
\documentclass{article}
\usepackage{colortbl}
\usepackage{makecell}
\usepackage{multirow}
\usepackage{supertabular}

\begin{document}

\newcounter{utterance}

\centering \large Interaction Transcript for game `wordle', experiment `medium\_frequency\_words\_no\_clue\_no\_critic', episode 9 with glm{-}5.2.
\vspace{24pt}

{ \footnotesize  \setcounter{utterance}{1}
\setlength{\tabcolsep}{0pt}
\begin{supertabular}{c@{$\;$}|p{.15\linewidth}@{}p{.15\linewidth}p{.15\linewidth}p{.15\linewidth}p{.15\linewidth}p{.15\linewidth}}
    \# & \multicolumn{2}{c}{Player} && \multicolumn{2}{c}{Game Master} \\
    \hline

    \theutterance \stepcounter{utterance}  
    & & & \multicolumn{4}{p{0.6\linewidth}}{
        \cellcolor[rgb]{0.9,0.9,0.9}{
            \makecell[{{p{\linewidth}}}]{
                \texttt{\tiny{[P1$\langle$GM]}}
                \texttt{Eres un mago del lenguaje al que le gusta adivinar palabras usando las reglas dadas.} \\
\\ 
\texttt{¡Bienvenido a Wordle! Tienes seis intentos para adivinar la palabra objetivo, una palabra válida en español de cinco letras minúsculas. Por favor, usa las etiquetas "explicación:" e "intento:" para proporcionar una explicación concisa para cada intento.} \\
\\ 
\texttt{Por ejemplo, si tu intento es "parte", tu respuesta debería ser} \\
\texttt{explicación: esta es una palabra común en español de cinco letras, y estoy comenzando mi intento con esta palabra.} \\
\texttt{intento: parte} \\
\\ 
\texttt{Después de cada intento, tu respuesta será validada y recibirás retroalimentación indicando qué letras son correctas (verde), qué letras son correctas pero están en la posición incorrecta (amarillo) y qué letras son incorrectas (rojo). Esta retroalimentación puede ser útil para determinar qué letras incluir o excluir en tu siguiente intento.} \\
\\ 
\texttt{Por ejemplo, la retroalimentación para "parte" podría ser:} \\
\texttt{retroalimentación\_intento: p<amarillo> a<amarillo> r<verde> t<amarillo> e<rojo>} \\
\\ 
\texttt{La explicación debe contener detalles sobre cómo se usa el retroalimentación\_intento para llegar a un nuevo intento.} \\
\\ 
\texttt{No generes ningún otro texto.} \\
\\ 
\texttt{Comencemos con tu primer intento.} \\
            }
        }
    }
     \\ \\

    \theutterance \stepcounter{utterance}  
    & \multicolumn{4}{p{0.6\linewidth}}{
        \cellcolor[rgb]{0.561,0.737,0.561}{
            \makecell[{{p{\linewidth}}}]{
                \texttt{\tiny{[P1$\rangle$GM]}}
                \texttt{explicación: elijo una palabra común en español de cinco letras que incluye vocales frecuentes y consonantes comunes para abrir el juego.} \\
\texttt{intento: patio} \\
            }
        }
    }
    & & \\ \\

    \theutterance \stepcounter{utterance}  
    & & & \multicolumn{2}{p{0.3\linewidth}}{
        \cellcolor[rgb]{0.9,0.9,0.9}{
            \makecell[{{p{\linewidth}}}]{
                \texttt{\tiny{[GM$|$GM]}}
                \texttt{attempts = 1} \\
\texttt{target\_word = ánimo} \\
\texttt{guess = patio} \\
\texttt{guess\_feedback = p<red> a<red> t<red> i<yellow> o<green>} \\
            }
        }
    }
    & & \\ \\

    \theutterance \stepcounter{utterance}  
    & & & \multicolumn{4}{p{0.6\linewidth}}{
        \cellcolor[rgb]{0.9,0.9,0.9}{
            \makecell[{{p{\linewidth}}}]{
                \texttt{\tiny{[P1$\langle$GM]}}
                \texttt{retroalimentación\_intento: p<rojo> a<rojo> t<rojo> i<amarillo> o<verde>} \\
\\ 
\texttt{Proporciona tu respuesta solo en este formato.} \\
\texttt{explicación: detalles} \\
\texttt{intento: palabra} \\
            }
        }
    }
     \\ \\

    \theutterance \stepcounter{utterance}  
    & \multicolumn{4}{p{0.6\linewidth}}{
        \cellcolor[rgb]{0.561,0.737,0.561}{
            \makecell[{{p{\linewidth}}}]{
                \texttt{\tiny{[P1$\rangle$GM]}}
                \texttt{explicación: La "o" es correcta y está al final. La "i" está en la palabra pero no en la cuarta posición. Se descartan "p", "a" y "t". Busco palabras con "i" y "o" final, sin las descartadas, como "radio" (colocando la "i" en tercera posición).} \\
\texttt{intento: radio} \\
            }
        }
    }
    & & \\ \\

    \theutterance \stepcounter{utterance}  
    & & & \multicolumn{2}{p{0.3\linewidth}}{
        \cellcolor[rgb]{0.9,0.9,0.9}{
            \makecell[{{p{\linewidth}}}]{
                \texttt{\tiny{[GM$|$GM]}}
                \texttt{attempts = 2} \\
\texttt{target\_word = ánimo} \\
\texttt{guess = radio} \\
\texttt{guess\_feedback = r<red> a<red> d<red> i<yellow> o<green>} \\
            }
        }
    }
    & & \\ \\

    \theutterance \stepcounter{utterance}  
    & & & \multicolumn{4}{p{0.6\linewidth}}{
        \cellcolor[rgb]{0.9,0.9,0.9}{
            \makecell[{{p{\linewidth}}}]{
                \texttt{\tiny{[P1$\langle$GM]}}
                \texttt{retroalimentación\_intento: r<rojo> a<rojo> d<rojo> i<amarillo> o<verde>} \\
\\ 
\texttt{Proporciona tu respuesta solo en este formato.} \\
\texttt{explicación: detalles} \\
\texttt{intento: palabra} \\
            }
        }
    }
     \\ \\

    \theutterance \stepcounter{utterance}  
    & \multicolumn{4}{p{0.6\linewidth}}{
        \cellcolor[rgb]{0.561,0.737,0.561}{
            \makecell[{{p{\linewidth}}}]{
                \texttt{\tiny{[P1$\rangle$GM]}}
                \texttt{explicación: La "o" sigue estando en el lugar correcto (quinta posición) y la "i" sigue en la palabra pero no va en la cuarta. Se descartan "r", "a" y "d". Necesito una palabra con "i" en otra posición (segunda o tercera) y "o" al final. Una palabra que encaja y usa letras no descartadas es "mírio" (relativo al mirio), aunque para ser más seguro con vocabulario común sin "r" exploro opciones como "novio".} \\
\texttt{intento: novio} \\
            }
        }
    }
    & & \\ \\

    \theutterance \stepcounter{utterance}  
    & & & \multicolumn{2}{p{0.3\linewidth}}{
        \cellcolor[rgb]{0.9,0.9,0.9}{
            \makecell[{{p{\linewidth}}}]{
                \texttt{\tiny{[GM$|$GM]}}
                \texttt{attempts = 3} \\
\texttt{target\_word = ánimo} \\
\texttt{guess = novio} \\
\texttt{guess\_feedback = n<yellow> o<red> v<red> i<yellow> o<green>} \\
            }
        }
    }
    & & \\ \\

    \theutterance \stepcounter{utterance}  
    & & & \multicolumn{4}{p{0.6\linewidth}}{
        \cellcolor[rgb]{0.9,0.9,0.9}{
            \makecell[{{p{\linewidth}}}]{
                \texttt{\tiny{[P1$\langle$GM]}}
                \texttt{retroalimentación\_intento: n<amarillo> o<rojo> v<rojo> i<amarillo> o<verde>} \\
\\ 
\texttt{Proporciona tu respuesta solo en este formato.} \\
\texttt{explicación: detalles} \\
\texttt{intento: palabra} \\
            }
        }
    }
     \\ \\

    \theutterance \stepcounter{utterance}  
    & \multicolumn{4}{p{0.6\linewidth}}{
        \cellcolor[rgb]{0.561,0.737,0.561}{
            \makecell[{{p{\linewidth}}}]{
                \texttt{\tiny{[P1$\rangle$GM]}}
                \texttt{explicación: La "o" final se confirma. La "i" sigue en amarillo (no está en la cuarta posición), así que debe ir en la segunda o tercera. La "n" está en la palabra pero no al principio, así que la muevo a la tercera posición. Se descartan "o" (al inicio) y "v". Pruebo "inneo".} \\
\texttt{intento: inneo} \\
            }
        }
    }
    & & \\ \\

    \theutterance \stepcounter{utterance}  
    & & & \multicolumn{2}{p{0.3\linewidth}}{
        \cellcolor[rgb]{0.9,0.9,0.9}{
            \makecell[{{p{\linewidth}}}]{
                \texttt{\tiny{[GM$|$GM]}}
                \texttt{Error: The guessed word is not a valid word for this game.} \\
            }
        }
    }
    & & \\ \\

    \theutterance \stepcounter{utterance}  
    & & & \multicolumn{4}{p{0.6\linewidth}}{
        \cellcolor[rgb]{0.9,0.9,0.9}{
            \makecell[{{p{\linewidth}}}]{
                \texttt{\tiny{[P1$\langle$GM]}}
                \texttt{Tu intento no es una palabra válida para este juego. Por favor, inténtalo de nuevo.} \\
\\ 
\texttt{Proporciona tu respuesta solo en este formato.} \\
\texttt{explicación: detalles} \\
\texttt{intento: palabra} \\
            }
        }
    }
     \\ \\

    \theutterance \stepcounter{utterance}  
    & \multicolumn{4}{p{0.6\linewidth}}{
        \cellcolor[rgb]{0.561,0.737,0.561}{
            \makecell[{{p{\linewidth}}}]{
                \texttt{\tiny{[P1$\rangle$GM]}}
                \texttt{explicación: El intento anterior no fue válido. Mantengo la "o" al final y la "i" en la segunda posición, y coloco la "n" en la primera posición, que es otra posibilidad válida. La palabra "ninfo" (relativo a ninfo) es un término aceptado en español.} \\
\texttt{intento: ninfo} \\
            }
        }
    }
    & & \\ \\

    \theutterance \stepcounter{utterance}  
    & & & \multicolumn{2}{p{0.3\linewidth}}{
        \cellcolor[rgb]{0.9,0.9,0.9}{
            \makecell[{{p{\linewidth}}}]{
                \texttt{\tiny{[GM$|$GM]}}
                \texttt{attempts = 4} \\
\texttt{target\_word = ánimo} \\
\texttt{guess = ninfo} \\
\texttt{guess\_feedback = n<yellow> i<yellow> n<red> f<red> o<green>} \\
            }
        }
    }
    & & \\ \\

    \theutterance \stepcounter{utterance}  
    & & & \multicolumn{4}{p{0.6\linewidth}}{
        \cellcolor[rgb]{0.9,0.9,0.9}{
            \makecell[{{p{\linewidth}}}]{
                \texttt{\tiny{[P1$\langle$GM]}}
                \texttt{retroalimentación\_intento: n<amarillo> i<amarillo> n<rojo> f<rojo> o<verde>} \\
\\ 
\texttt{Proporciona tu respuesta solo en este formato.} \\
\texttt{explicación: detalles} \\
\texttt{intento: palabra} \\
            }
        }
    }
     \\ \\

    \theutterance \stepcounter{utterance}  
    & \multicolumn{4}{p{0.6\linewidth}}{
        \cellcolor[rgb]{0.561,0.737,0.561}{
            \makecell[{{p{\linewidth}}}]{
                \texttt{\tiny{[P1$\rangle$GM]}}
                \texttt{explicación: La "o" final se confirma. La "n" está en la palabra pero no al inicio; la "i" tampoco va en la segunda posición. Se descartan "f" y la segunda "n". Muevo la "i" a la tercera posición y la "n" a la cuarta, formando una palabra válida.} \\
\texttt{intento: dorio} \\
            }
        }
    }
    & & \\ \\

    \theutterance \stepcounter{utterance}  
    & & & \multicolumn{2}{p{0.3\linewidth}}{
        \cellcolor[rgb]{0.9,0.9,0.9}{
            \makecell[{{p{\linewidth}}}]{
                \texttt{\tiny{[GM$|$GM]}}
                \texttt{attempts = 5} \\
\texttt{target\_word = ánimo} \\
\texttt{guess = dorio} \\
\texttt{guess\_feedback = d<red> o<red> r<red> i<yellow> o<green>} \\
            }
        }
    }
    & & \\ \\

    \theutterance \stepcounter{utterance}  
    & & & \multicolumn{4}{p{0.6\linewidth}}{
        \cellcolor[rgb]{0.9,0.9,0.9}{
            \makecell[{{p{\linewidth}}}]{
                \texttt{\tiny{[P1$\langle$GM]}}
                \texttt{retroalimentación\_intento: d<rojo> o<rojo> r<rojo> i<amarillo> o<verde>} \\
\\ 
\texttt{Proporciona tu respuesta solo en este formato.} \\
\texttt{explicación: detalles} \\
\texttt{intento: palabra} \\
            }
        }
    }
     \\ \\

    \theutterance \stepcounter{utterance}  
    & \multicolumn{4}{p{0.6\linewidth}}{
        \cellcolor[rgb]{0.561,0.737,0.561}{
            \makecell[{{p{\linewidth}}}]{
                \texttt{\tiny{[P1$\rangle$GM]}}
                \texttt{explicación: La "o" final sigue correcta y la "i" sigue en amarillo, lo que indica que debe ir en la segunda o tercera posición. Se descartan "d", "o" (al inicio) y "r". Busco una palabra con "n" (que aún no ha sido fijada), "i" en segunda o tercera, y "o" final. La palabra "chico" encaja, pero la "c" aún no se ha probado. Sin embargo, recordando que "n" estaba en amarillo antes, intento colocar "n" al inicio y "i" en tercera posición.} \\
\texttt{intento:Union} \\
            }
        }
    }
    & & \\ \\

    \theutterance \stepcounter{utterance}  
    & & & \multicolumn{2}{p{0.3\linewidth}}{
        \cellcolor[rgb]{0.9,0.9,0.9}{
            \makecell[{{p{\linewidth}}}]{
                \texttt{\tiny{[GM$|$GM]}}
                \texttt{attempts = 6} \\
\texttt{target\_word = ánimo} \\
\texttt{guess = union} \\
\texttt{guess\_feedback = u<red> n<green> i<green> o<yellow> n<red>} \\
            }
        }
    }
    & & \\ \\

    \theutterance \stepcounter{utterance}  
    & & & \multicolumn{2}{p{0.3\linewidth}}{
        \cellcolor[rgb]{0.9,0.9,0.9}{
            \makecell[{{p{\linewidth}}}]{
                \texttt{\tiny{[GM$|$GM]}}
                \texttt{game\_result = LOSS} \\
            }
        }
    }
    & & \\ \\

\end{supertabular}
}

\end{document}
